\chapter{Developing the doctrinal analysis}\label{ch:doctrine}
\section{Move from the source to the proposition}
Identify the proposition, the source's status and the question it addressed.
Distinguish a judgment's reasoning from other observations. Read legislation
with its definitions, exceptions and related provisions. Explain how the text
and context support your conclusion.

A first reference in this chapter is full again.\lawcite[{[16]}]{vellama2013}
A repeat points to this chapter's first note.
\lawcite[{[16]}]{vellama2013}
The table of cases collects the cited thesis pages.

A European Union case can be recorded with a different jurisdiction and
its own court identifier.\lawcite[{[8]}]{marleasing1990}
It appears under a separate jurisdiction heading in the table of cases.

\section{Version-sensitive legislation}
The source register identifies the Human Rights Act example as the original,
as-enacted provision.\lawcite[s~3(1)]{hra1998}
That choice demonstrates how to record a statutory version. A thesis making a
claim about current law must examine the version and legal developments relevant
to its own date and jurisdiction.

\section{Footnotes that contain analysis}
A note can explain a source choice, cite more than one authority and retain the
same citation tracking.\footnote{For the legislative example, see
\lawref[s~3(1)]{hra1998}. For the treaty example, see
\lawref[art~31(1)]{vclt1969}. The examples illustrate two different source types.}
Use this form where the relationship between authorities needs explanation.
Place extended reasoning in the main text when the argument depends on it.

\section{Counter-authority and the strongest objection}
Identify the strongest counter-authority and state its force fairly.
Explain whether jurisdiction, hierarchy, date, facts or legal principle
justifies a distinction. If the objection survives, narrow the claim.

\LawPrompt{For each principal proposition, identify its best supporting
authority, its best counter-authority and the exact issue that separates them.}
